\section{Introdução}

A mobilidade ativa, isto é, a mobilidade cujos deslocamentos são realizados a pé, de bicicleta ou por algum outro meio que não envolva a utilização de veículos individuais motorizados, pode contribuir na mitigação de impactos climáticos. Segundo \cite{BRAND2021102764}, cerca de 70\% das emissões de ciclo de vida de CO2 por pessoa de atividades relacionadas a mobilidade urbana são geradas por carros frente a apenas 1\% das emissões geradas por bicicleta. Os autores ainda estimam que cada deslocamento de carro evitado reduz em 62\% as emissões por indivíduo.

Em resposta à urgência climática e em busca da criação de políticas de descarbonização, diversos modelos de planejamento urbano vêm sendo desenvolvidos. Um modelo de destaque, atualmente em implantação em importantes cidades como Paris e Barcelona, é o modelo de cidade de 15 minutos \cite{Allam2022-yx}. Neste modelo, a cidade é organizada de forma que os principais serviços utilizados no dia a dia da população possam ser acessados através de deslocamentos de curta distância, priorizando a mobilidade ativa.

Contudo, a viabilidade de modelos como a cidade de 15 minutos depende fundamentalmente de como a rede de deslocamentos se comporta sob condições reais de operação. Trabalhos anteriores demonstraram que a acessibilidade urbana não é uniforme: desigualdades socioeconômicas moldam fortemente quem tem acesso a quais serviços e em qual tempo \cite{sbbd_estendido}. Além disso, eventos climáticos extremos (inundações, chuvas intensas, escorregamentos) degradam a infraestrutura de mobilidade, interrompendo ou dificultando trechos da rede viária e tornando certos deslocamentos inviáveis ou significativamente mais custosos. Essas perturbações não afetam a população de forma uniforme: áreas mais vulneráveis frequentemente apresentam menor resiliência à degradação da rede e populações com menor capacidade de adaptação sofrem maiores impactos.

Este trabalho avalia como perturbações climáticas modificam a acessibilidade urbana e como esses impactos se distribuem desigualmente no espaço e entre perfis socioeconômicos, propondo uma abordagem metodológica que integra modelagem de redes de deslocamento com análise socioeconômica e compara cenários típicos (sem perturbação) e cenários condicionados por variáveis ambientais. Com isso, contribui para a compreensão de como garantir equidade na acessibilidade urbana sob cenários de perturbação ambiental.